\begin{figure}[!htbp]
\begin{examplepanel}{srccol}{Source function subtyping --- contravariant domain}
\[
\src{\inferrule*[right=\textsc{All}]{
 W<:\top\\ \top<:\top
}{(\top\to\top)<:(W\to\top)}}
\]
\end{examplepanel}
\begin{examplepanel}{tgtcol}{Target evidence and application}
\[
\tgt{\begin{aligned}
 d&=\ev{top}(W):W\le\top,&c&=\ev{refl}(\top):\top\le\top,\\
 g&=\cast{z}{\ev{pi}(d,c)}:W\to\top.
\end{aligned}}
\]
\par\small
\renewcommand{\arraystretch}{1.2}
\noindent\begin{tabularx}{\linewidth}{@{}lX@{}}
\toprule
Stage of \(g\,z\) & Argument or result at that stage\\
\midrule
Client supplies & $z:W$\\
Original function receives & $\cast{z}{d}:\top$\\
Identity body returns & $\cast{z}{d}:\top$\\
Client receives & $\cast{\cast{z}{d}}{c}:\top$\\
\bottomrule
\end{tabularx}\par
\end{examplepanel}
\begin{examplepanel}{runcol}{Erased execution}
\[
\run{\erase{\sigma}(z)=\lambda u.u,\qquad z\,z\longrightarrow z.}
\]
\end{examplepanel}
\caption{The function coercion adapts the argument with \(d\) and the
result with \(c\). The typed store holds \(z=\lambda(u:\top).u\), and
the function head form retains both evidence components for application.
The table groups administrative steps; the erased call applies the
stored identity.}
\label{fig:function-call}
\end{figure}
